% Einheit 3 von 3 – Mit Parameter (bank/bedingte-wahrscheinlichkeit-und-bayes/e3.jsonl, zusammenbau v0.1)
%% TODO Zeitmarke Einheit 3: Marken-Zeile nicht lesbar
%% TODO Zweigzeile Teil 1: was hier gelernt wird (Satz in Schülersprache aus dem Einheitstitel) – Einheit 3
\einheitenkopf[e3]{Einheit 3 von 3 · Mit Parameter}
\zweigzeile{Abitur GK}
\setcounter{aufgabe}{11}

%% TODO Titel: Ich-kann-Satz fehlt in der Bank (Platzhalter: Kettenname) – Nr. 12
%% TODO Anweisung: ein Satz über der ersten Teilaufgabe fehlt in der Bank – Nr. 12
\begin{aufgabe}{Mit Parameter}
%% TODO \rechenplatz{n}: Schrittzahl des Grundfalls fehlt in der Bank (nur bei fester Schreibform, 2.3 b)
\begin{teile}
\teil Wo steht der Parameter im Term $\frac{0{,}3}{0{,}3 + a}$? \\ \kreuz{nur im Zähler} \\ \kreuz{nur im Nenner} \\ \kreuz{in Zähler und Nenner}
\teil Ein Anteil a der Kunden eines Geschäfts, $0 < a < 1$, sind Stammkunden. 60\,\% der Stammkunden und 10\,\% der übrigen Kunden kaufen ein Angebot. Ein Kunde kauft das Angebot. Stelle einen Term in a für die Wahrscheinlichkeit auf, dass er Stammkunde ist. \leerfeld
\teil Ein Drittel der Kunden eines Versandhändlers sind Neukunden. 10\,\% der Neukunden und ein Anteil b der Stammkunden reklamieren. Durch bessere Beratung sinkt b. Begründe ohne Rechnung, dass dabei die Wahrscheinlichkeit wächst, dass ein reklamierender Kunde ein Neukunde ist.
\teil 30\,\% der Besucher eines Freibads kommen mit dem Rad; von ihnen kauft die Hälfte ein Eis, von den übrigen Besuchern ein Anteil a. Ein Besucher kauft ein Eis. Welcher der Graphen A, B, C zeigt die Wahrscheinlichkeit, dass er mit dem Rad kam, in Abhängigkeit von a? Begründe ohne Rechnung über die Randwerte.

\begin{ksys}[xmin=0,xmax=1,ymin=0,ymax=1,xstep=0.2,ystep=0.2,ablesen] \funktion{0.15/(0.15+0.7*\x)}{A} \funktion{1-\x}{B} \funktion{0.15/(0.15+0.7*(1-\x))}{C} \end{ksys}
\teil 40\,\% der Gäste eines Hotels reisen geschäftlich; 75\,\% von ihnen nutzen das WLAN, von den übrigen Gästen ein Anteil a. Der Graph zeigt $f(a)$, die Wahrscheinlichkeit, dass ein WLAN-Nutzer geschäftlich reist. Bestimme am Graphen a mit $f(a) = 0{,}5$ und deute a und $f(a)$ im Sachzusammenhang.

\begin{ksys}[xmin=0,xmax=1,ymin=0,ymax=1,xstep=0.1,ystep=0.1,ablesen] \funktion{0.3/(0.3+0.6*\x)}{f} \end{ksys}
\teil Bei einer Apfelernte ist jeder zehnte Apfel faulig. Eine Sortiermaschine stuft jeden guten Apfel als gut ein und erkennt einen fauligen mit der Wahrscheinlichkeit x als faulig. Ein als gut eingestufter Apfel wird zufällig ausgewählt. Wie groß muss x mindestens sein, damit er höchstens mit 2\,\% faulig ist? \hfill (Abitur 2018 LK)
\end{teile}
\end{aufgabe}

%% TODO Titel: Ich-kann-Satz fehlt in der Bank (Platzhalter: Kettenname) – Nr. 13
%% TODO Anweisung: ein Satz über der ersten Teilaufgabe fehlt in der Bank – Nr. 13
\begin{aufgabe}{Mit Parameter – Fehler finden, Begründen, Anwendung}
\begin{teile}
\teil Lisa soll begründen, dass $\frac{0{,}12}{0{,}12 + 0{,}5b}$ wächst, wenn b fällt. Sie schreibt: „Für $b = 0{,}4$ kommt $0{,}375$ heraus, für $b = 0{,}2$ kommt $0{,}545$ heraus, also wächst der Bruch.“ Finde den Fehler.
\teil Begründe, warum ein Bruch mit positivem, konstantem Zähler wächst, wenn sein positiver Nenner fällt.
\teil 1\,\% der Bevölkerung hat eine bestimmte Krankheit. Ein Test ist bei Kranken mit 95\,\% positiv; bei Gesunden ist er mit der Wahrscheinlichkeit w falsch positiv. Wie klein muss w höchstens sein, damit eine positiv getestete Person mindestens mit 50\,\% wirklich krank ist?
\end{teile}
\end{aufgabe}
